% ===================================================================
\section{Novel Results from Formalization}
\label{ch:results}
% ===================================================================

\subsection{Overview}

Three non-trivial theorems emerge directly from the Lean formalization --- results
not present in the original published theory and arising solely from the proof
obligations imposed by the formal system. This chapter presents these theorems,
their proofs, their interpretations, and the open questions they suggest.

Additionally, 8 direct consequence theorems from the reduced basis are catalogued.

% -------------------------------------------------------------------
\subsection{Novel Theorem 1: First Unresolved Witness}
% -------------------------------------------------------------------

\begin{theorem}[First unresolved witness]
If an unresolved ambiguity exists anywhere in a protocol trace, there is a
\emph{canonical first} such occurrence:
\begin{align*}
  &\forall \mathcal{T},\
  \Big(\exists n,\
  \mathcal{T}[n].\mathrm{ambiguous} = \top
  \wedge \mathcal{T}[n].\mathrm{clarified} = \bot\Big)
  \Rightarrow \\
  &\exists n^*,\
  \Big(\mathcal{T}[n^*].\mathrm{ambiguous} = \top
  \wedge \mathcal{T}[n^*].\mathrm{clarified} = \bot\Big)
  \wedge \\
  &\quad \Big(\forall m < n^*,\
  \neg(\mathcal{T}[m].\mathrm{ambiguous} = \top
      \wedge \mathcal{T}[m].\mathrm{clarified} = \bot)\Big)
\end{align*}
\leanref{Frfp.Minimal.BasisSweep.first\_unresolved\_witness}
--- machine-checked proof.
\end{theorem}
\begin{proof}
Let $P(n) = \text{"step } n \text{ is unresolved"}$. By hypothesis, $\exists n,
P(n)$. By Lean's \texttt{Nat.find}, there exists a minimal witness $n^* =
\mathrm{Nat.find}(P)$ satisfying $P(n^*)$, and for all $m < n^*$, $\neg P(m)$.
\qed
\end{proof}

\paragraph{Why novel.} The original AR axiom says "all ambiguities are resolved"
--- a global property. This theorem localizes AR violations: when AR fails, it
fails at a specific first step. This is essential for:
\begin{itemize}
  \item \textbf{Audit trails}: governance audits can find the exact first violation
        point.
  \item \textbf{Incremental checking}: a trace is AR-compliant up to step $n$ iff
        $n < n^*$.
  \item \textbf{Repair protocols}: a minimal repair targets only step $n^*$, not
        the entire trace.
\end{itemize}

% -------------------------------------------------------------------
\subsection{Novel Theorem 2: Governance Safe Extension}
% -------------------------------------------------------------------

\begin{theorem}[Governance safe extension]
Strong governance admissibility (IL, AR, CSC) is preserved under compatible
one-step trace extension:
\[
  \mathrm{GA}(\mathcal{T}) \wedge \mathrm{Compatible}(\mathcal{T}, s)
  \Rightarrow
  \mathrm{GA}(\mathcal{T} \mathbin{+\!\!+} [s])
\]
\leanref{Frfp.Minimal.BasisSweep.governance\_safe\_extension\_theorem}
--- machine-checked proof.
\end{theorem}
\begin{proof}
Each of the three conjuncts of $\mathrm{GA}$ is preserved independently:
\begin{description}
  \item[\textbf{IL}:] For any two AI steps $a, b$ in $\mathcal{T}
        \mathbin{+\!\!+} [s]$, if both are in $\mathcal{T}$ then
        $\mathrm{IL}_{\mathrm{global}}(\mathcal{T})$ applies; if one is $s$ and
        the other is in $\mathcal{T}$, the intent compatibility condition applies.
  \item[\textbf{AR}:] For any ambiguous step in $\mathcal{T}
        \mathbin{+\!\!+} [s]$: if in $\mathcal{T}$,
        $\mathrm{AR}_{\mathrm{global}}(\mathcal{T})$ applies; if it is $s$, the
        ambiguity compatibility condition applies.
  \item[\textbf{CSC}:] The lock-forward and lock-backward conditions of
        compatibility ensure the lock propagates consistently. \qed
\end{description}
\end{proof}

\paragraph{Why novel.} The original FRFP governance axioms are stated as global
properties of a complete trace. Theorem 7.2 is an \emph{inductive} property ---
it says governance compliance is closed under one-step extension. This enables:
\begin{itemize}
  \item \textbf{Online governance checking}: evaluate each new step independently.
  \item \textbf{Compositional reasoning}: compliant sub-traces compose to a
        compliant whole.
  \item \textbf{Distributed protocols}: each agent can enforce local step
        compatibility; global compliance follows.
\end{itemize}

% -------------------------------------------------------------------
\subsection{Novel Theorem 3: Protocol Normal Form}
% -------------------------------------------------------------------

\begin{theorem}[Protocol normal form]
Every strongly admissible primitive trace decomposes into an explicit block, an
optional boundary crossing, and a tacit block:
\begin{align*}
  &\forall \mathcal{T} : \mathrm{List}(\mathrm{Primitive}),\
  \mathrm{AdmissibleStrong}(\mathcal{T}) \Rightarrow \\
  &\exists\ e_1\ldots e_k,\ h_1\ldots h_m,\ b \in \mathbb{B}, \\
  &\mathcal{T}
  = [e_1,\ldots,e_k]
    \mathbin{+\!\!+}
    (\text{if } b \text{ then } [\RB] \text{ else } [])
    \mathbin{+\!\!+}
    [h_1,\ldots,h_m]
\end{align*}
where each $e_i \in \Ecat = \{\RI, \EC, \ED\}$ and each
$h_j \in \Tcat \setminus \{\RB\} = \{\TE, \HFD\}$.

\leanref{Frfp.Minimal.BasisSweep.protocol\_normal\_form\_theorem}
--- machine-checked proof.
\end{theorem}

\paragraph{Why novel.} The three-block structure is mentioned informally in the
original theory. Theorem 7.3 makes it a formal theorem:
\begin{itemize}
  \item \textbf{Every} admissible trace has this structure --- it is not optional.
  \item The decomposition is canonical: the explicit block always precedes $\RB$,
        which always precedes the tacit block.
  \item The boundary $\RB$ appears \textbf{at most once} (by Theorem 1.2).
\end{itemize}

\paragraph{Implication for protocol design.} Any protocol that claims to be
FRFP-compliant but allows $\RB$ in a non-terminal position, or that interleaves
explicit and tacit steps, is not strongly admissible.

% -------------------------------------------------------------------
\subsection{Direct Consequences of the Reduced Basis}
% -------------------------------------------------------------------

\begin{theorem}[Tacit source classification]
$\forall p : \mathrm{Primitive},\ p.\mathrm{source} = \Tobj
\Rightarrow p \in \{\TE, \HFD\}$
\leanref{Frfp.Minimal.BasisConsequences.tacit\_source\_classification}
\end{theorem}

\begin{theorem}[Tacit source stays tacit]
$\forall p,\ p.\mathrm{source} = \Tobj \Rightarrow p.\mathrm{target} = \Tobj$
\leanref{Frfp.Minimal.BasisConsequences.tacit\_source\_stays\_tacit}
\end{theorem}

\begin{theorem}[AI steps never target tacit space]
$\forall p,\ \textit{is\_AI\_executable}(p) = \top \Rightarrow p.\mathrm{target} \neq \Tobj$
\leanref{Frfp.Minimal.BasisConsequences.ai\_never\_targets\_tacit}
\end{theorem}

\begin{theorem}[AI/human disjoint]
$\forall p,\ \neg(\textit{is\_AI\_executable}(p) \wedge \textit{is\_human\_only}(p))$
\leanref{Frfp.Minimal.BasisConsequences.ai\_human\_disjoint}
\end{theorem}

\begin{theorem}[Unique boundary entry]
$|\{p : \mathrm{Primitive} \mid p.\mathrm{source} = \Eobj \wedge p.\mathrm{target} = \Tobj\}| = 1$
\leanref{Frfp.Minimal.BasisConsequences.unique\_boundary\_entry}
\end{theorem}

\begin{theorem}[IL pairwise intent consistency]
Under IL, any two AI steps in the same trace share an intent tag.
\leanref{Frfp.Minimal.BasisConsequences.IL\_pairwise\_intent\_consistency}
\end{theorem}

\begin{theorem}[AR: no unresolved ambiguity]
Under AR, no ambiguous step remains unclarified.
\leanref{Frfp.Minimal.BasisConsequences.AR\_no\_unresolved\_ambiguity}
\end{theorem}

\begin{theorem}[CSC: lock is monotone]
Under CSC, the correctness-locked flag is monotone non-decreasing over the trace.
\leanref{Frfp.Minimal.BasisConsequences.CSC\_lock\_is\_monotone}
\end{theorem}

% -------------------------------------------------------------------
\subsection{Novelty Classification}
% -------------------------------------------------------------------

\begin{center}
\begin{tabular}{lll}
\toprule
Theorem & Type & How it arose \\
\midrule
7.1 First unresolved witness & New result & \texttt{Nat.find} forced minimal witness existence \\
7.2 Governance safe extension & New inductive property & Needed to prove $\mathrm{GA}$ is inductive \\
7.3 Protocol normal form & New completeness theorem & Admissibility definition crystallized three-block shape \\
Thms 7.4--7.11 & Direct consequences & Short proofs from basis definitions \\
\bottomrule
\end{tabular}
\end{center}

% -------------------------------------------------------------------
\subsection{Open Questions}
% -------------------------------------------------------------------

\begin{enumerate}[label=\textbf{Q\arabic*}]
  \item \textbf{Complete $\runToOmega$ proof.} The injectivity and round-trip
        properties are axiomatized. A full proof requires constructing a
        measurable bijection between the discrete run type and the
        measure-theoretic trajectory space.

  \item \textbf{Phase 2 primitives.} The current kernel is Phase 1 (6
        primitives). Can a Phase 2 extension be defined such that the analogues
        of ETS and RB uniqueness remain derivable theorems?

  \item \textbf{Mechanized minimality.} The minimality of the 7-axiom basis
        is established by manual witness construction. An automated model-finder
        could confirm or refute this mechanically.

  \item \textbf{Semantic grounding decidability.} The $\mathrm{CorrectnessJudgment}$
        type has value $\mathrm{unknown}$ for undecidable cases. Is there a
        formal characterization of when grounding is decidable vs.\ undecidable?

  \item \textbf{Probabilistic confluence.} Theorems 5.3--5.7 are deterministic
        properties of the hazard/survival functions. Is there a probabilistic
        confluence theorem: does the distribution over normal forms converge
        under repeated random reduction?
\end{enumerate}
